% p2-06-phi-operator

\section{P2 §6. $\varphi$-Structured Operator Framework}

6. $\varphi$-Structured Operator Framework

           Figure (fig-p2-ch6-phi-operator). phi-operator triptych --- Operator action / phi-tower / 4D
                                             structural obstruction.

  [CONJECTURAL --- L5 throughout]

  6.1 Candidate Operator Setup
  Consider a 2D CFT with central charge c and a set of primary operators On with scaling
  dimensions Deltan. Define a $\varphi$-tower as a sequence of operators with

  Deltan = Delta0 $\cdot$ $\varphi$n, n = 0, 1, 2, … 16

  For such a tower to arise from a consistent CFT, the dimensions Deltan must satisfy Kac
  determinant positivity conditions and OPE closure. In a Virasoro minimal model M(p,p')
  with c < 1, scaling dimensions are rational numbers

  Deltar,s = ((rp' - sp)2 - (p'-p)2)/(4pp'), 17

  for integers r, s and the pair (p, p') characterizing the model. A $\varphi$-tower would require
  Deltar',s'/Deltar,s = $\varphi$n, placing an algebraic constraint on (r,s,p,p') that must be
  checked for all models with p, p' $\leq$ 30.

  Existence question [WP1, open]. Is there a unitary CFT containing two primary
  operators with dimension ratio $\varphi$? A non-exhaustive search for minimal models with this
  property has not appeared in the published literature. This constitutes a well-defined
  mathematical question answerable without new experiments. It is Work Package 1 in
  Section 9, with the formal success/failure criterion given in Section 4.4.

  6.2 Perturbing Away from the $\varphi$-Tower: RG Interpretation [L5]
  Suppose such a tower exists in a 2D CFT. Consider a perturbation by the relevant
  operator O0:

  A = ACFT + lambda0 $\in$t d2 x O0(x). 18

  The OPE O0 $\times$ O0 will in general contain O0 and other operators in the tower,
  generating a flow in coupling space. The $\varphi$-tower conjecture asserts that this flow has
  an approximate invariance: the ratio of couplings lambdan / lambda0 is attracted
  toward $\varphi$-nalpha for some exponent alpha determined by the anomalous dimension.
  This is entirely conjectural. No such flow has been computed.

  6.3 Extension to 4D: Structural Obstacles [L4--L5]
  In 4D, the space of scaling dimensions is not constrained by a Kac formula but by the
  operator spectrum of the specific theory. Whether any 4D QFT fixed point has two
  operator dimensions in ratio $\varphi$ is open. Three structural obstacles deserve explicit
  treatment:

  1. Rationality. In 2D CFTs, scaling dimensions at rational points of moduli space are
  often algebraic. In 4D, loop corrections introduce transcendental numbers (zeta(3),
  etc.) into anomalous dimensions, making algebraic ratios unlikely to persist beyond
  fixed perturbative order.

  2. Non-renormalization. In supersymmetric theories, non-renormalization theorems
  protect certain operator dimensions exactly. A SUSY theory with a $\varphi$-eigenvalue ratio at
  tree level would maintain it to all orders --- providing a structural motivation to search
  within the landscape of supersymmetric CFTs (SCFTs).

  3. Large-N. In large-N limits, anomalous dimensions simplify and in some cases become
  rational. Whether $\varphi$-ratios arise in this limit is an open computational question.

  These obstacles do not constitute falsifications of the $\varphi$-tower conjecture; they identify
  the theoretical spaces in which the conjecture is most or least likely to be realized.

\subsection{References}

- Coldea, R. et al. "Quantum Criticality in an Ising Chain: Experimental Evidence for Emergent E8 Symmetry." \textit{Science} 327 (2010) 177. DOI \href{https://doi.org/10.1126/science.1180085}{10.1126/science.1180085}.
- Wu, J. et al. "E8 Symmetry in CoN$b_{2}$$O_{6}$ at Higher Resolution." \textit{Phys. Rev. B} 108, L060403 (2023). DOI \href{https://doi.org/10.1103/PhysRevB.108.L060403}{10.1103/PhysRevB.108.L060403}.
- Fring, A.; Korff, C. "Affine Toda field theories related to Coxeter groups of non-crystallographic type." \textit{Nucl. Phys. B} 729 (2005) 361. \href{https://arxiv.org/abs/hep-th/0506226}{arXiv:hep-th/0506226}. DOI \href{https://doi.org/10.1016/j.nuclphysb.2005.08.044}{10.1016/j.nuclphysb.2005.08.044}.
- Pellis, S. "The Fine-Structure Constant and the Golden Ratio." \href{https://vixra.org/pdf/2110.0117v4.pdf}{viXra 2110.0117 v4}.
